\section{Conclusion}

Token-level credit assignment and parameter-efficient fine-tuning are usually
studied as separate design choices. This paper argues that they are coupled.
Under LoRA, output-distribution signals such as surprisal, entropy reduction,
and policy divergence can degenerate into token weights that no longer track
where the adapter acts. ARCA addresses this failure mode by measuring
salience where LoRA actually acts: in the adapter-induced hidden-state
residual. The resulting signal is lightweight, requires no learned critic or
learned process reward model, and produces non-degenerate token-credit distributions in
the regime where output-distribution baselines become either flat or highly
concentrated. Our MATH/Qwen3-1.7B sweep is intentionally compact, but it
supports the central mechanism: adaptation-aware credit assignment changes the
geometry of token updates under matched rollout and parameter budgets.
